\documentclass[border=10pt]{standalone}
\usepackage{fontspec}
\setmainfont{Noto Serif}
\newfontfamily\bengalifont{Noto Serif Bengali}[Script=Bengali]
\usepackage[Latin,Bengali]{ucharclasses}
\setTransitionsFor{Bengali}{\begingroup\bengalifont}{\endgroup}
\usepackage{newunicodechar}
\newunicodechar{।}{{\bengalifont\char"0964}}
\newunicodechar{॥}{{\bengalifont\char"0965}}

\usepackage[table]{xcolor}
\usepackage{booktabs,tabularx,array}
\definecolor{bapilrule}{HTML}{1F2937}
\definecolor{bapilstripe}{HTML}{F3F4F6}
\renewcommand{\arraystretch}{1.28}
\setlength{\tabcolsep}{5pt}

\begin{document}
\begin{minipage}{17cm}
{\large\bfseries সারণি ৬৭।}\ {\large ঝুঁকি পর্যবেক্ষণ: আগাম সতর্কসংকেত, বিকল্প ব্যবস্থা ও অবশিষ্ট মাত্রা}\par\vspace{7pt}
\small
\rowcolors{2}{bapilstripe}{white}
\begin{tabularx}{\linewidth}{@{}>{\hsize=0.326\hsize\raggedright\arraybackslash}X>{\hsize=1.630\hsize\raggedright\arraybackslash}X>{\hsize=1.522\hsize\raggedright\arraybackslash}X>{\hsize=0.652\hsize\raggedright\arraybackslash}X>{\hsize=0.870\hsize\raggedright\arraybackslash}X@{}}
\arrayrulecolor{bapilrule}\toprule
\rowcolor{white}\textbf{আইডি} & \textbf{আগাম সতর্কসংকেত} & \textbf{ঝুঁকি ঘটলে বিকল্প ব্যবস্থা} & \textbf{অবশিষ্ট (স/\allowbreak{}প্র)} & \textbf{পর্যালোচনার পুনরাবৃত্তি} \\
\midrule
R1 & অনুরোধের ৯০ দিনের মধ্যে কোনো সরকারি ফোকাল পয়েন্ট মনোনীত না হওয়া; জ্যেষ্ঠ পৃষ্ঠপোষকের পরিবর্তন & উন্মুক্ত ডেটা নিয়ে বিশ্ববিদ্যালয়ের গবেষণা ল্যাব হিসেবে চালিয়ে যাওয়া; অন্য পৃষ্ঠপোষক খোঁজা & মা/\allowbreak{}মা & ত্রৈমাসিক \\
R2 & ৬০ দিনের মধ্যে ডেটা-রক্ষকের সাড়া না পাওয়া; একসঙ্গে প্রবেশাধিকার প্রত্যাখ্যাত & ইতিমধ্যে অনলাইনে থাকা নথি ব্যবহার; তথ্য অধিকার আবেদন; কর্পাস সীমিত করা & মা/\allowbreak{}মা & প্রথম বছরে মাসিক \\
R3 & ২০\%-এর বেশি পৃষ্ঠায় OCR মান-স্কোর ০.৯-এর নিচে & হাতে সংশোধন বাড়ানো; ডিজিটালভাবে তৈরি উৎসকে অগ্রাধিকার & মা/\allowbreak{}নি & প্রতিটি কর্পাস রিলিজে \\
R4 & উদ্ধৃতির নির্ভুলতা ০.৯৫-এর নিচে; ভুল উত্তরের প্রতিবেদন & শুধু অনুসন্ধান-মোডে ফেরা; ব্যবহারকারীদের পুনঃপ্রশিক্ষণ & নি/\allowbreak{}উ & প্রতিটি রিলিজে; পাইলটে মাসিক \\
R5 & লগে কোনো P2 ডেটা পাওয়া; প্রবেশাধিকারে অস্বাভাবিকতা & P3 স্থগিত; আইন অনুযায়ী অবহিতকরণ; ERRC পর্যালোচনা & নি/\allowbreak{}উ & মাসিক \\
R6 & ব্যক্তি-নজরদারির অনুরোধ; T0 স্তর-সংশ্লিষ্ট পরিধি-পরিবর্তন & ERRC-এর প্রত্যাখ্যান; সিদ্ধান্ত প্রকাশ & নি/\allowbreak{}উ & প্রতিটি গেটে \\
R7 & ৩ মাসের বেশি পদ শূন্য; ৩০\%-এর বেশি নিয়োগপ্রস্তাব প্রত্যাখ্যাত & ফেলো ও চুক্তিভিত্তিক কর্মী; দূরবর্তী কর্মী & মা/\allowbreak{}মা & ত্রৈমাসিক \\
R8 & মাস ১৮-এর মধ্যে দ্বিতীয় ধাপের অর্থায়ন-প্রতিশ্রুতি নেই & মূল কাজে সংকুচিত করা: কর্পাস, বেঞ্চমার্ক ও শাসন & মা/\allowbreak{}উ & ত্রৈমাসিক \\
R9 & এআই হাব প্রকল্পের পরিকল্পনায় সাংঘর্ষিক কার্যক্রম & যৌথ কর্মপরিকল্পনা; কাজের ধরন অনুযায়ী বিভাজন & নি/\allowbreak{}মা & প্রতি ছয় মাসে \\
R10 & বিকল্প-পরিকল্পনা ছাড়া একাধিক উপাদানে একক বিক্রেতা বা মডেল & পুনঃদরপত্র; বেঞ্চমার্কের ভিত্তিতে মডেল বদল & নি/\allowbreak{}মা & প্রতিটি ADR-এ \\
R11 & অনুপ্রবেশ পরীক্ষায় গুরুতর ত্রুটি; অনুপ্রবেশের সতর্কসংকেত & বিচ্ছিন্ন করা; অফলাইন ব্যাকআপ থেকে পুনরুদ্ধার & নি/\allowbreak{}উ & ত্রৈমাসিক পরীক্ষা; নিরবচ্ছিন্ন পর্যবেক্ষণ \\
R12 & টানা দুই মাস মাসিক সক্রিয় ব্যবহারকারী লক্ষ্যের অর্ধেকের নিচে & সহ-নকশা স্প্রিন্ট; পাইলট পুনর্বিবেচনা & মা/\allowbreak{}মা & মাসিক \\
R13 & নতুন বিধিমালা জারি; আইনি পর্যালোচনার ফলাফল & প্রক্রিয়া হালনাগাদ; প্রভাবিত প্রক্রিয়াকরণ স্থগিত & নি/\allowbreak{}মা & বিধিমালা জারির সময়; ত্রৈমাসিক \\
R14 & অপব্যবহারের প্রতিবেদন; সমষ্টিগত প্রশ্নধারায় অস্বাভাবিকতা & প্রবেশাধিকার সীমিত; ERRC পর্যালোচনা & নি/\allowbreak{}উ & মাসিক \\
R15 & উন্নয়ন ও পরীক্ষার ফলাফলে ৫ পয়েন্টের বেশি ব্যবধান; পরীক্ষা-প্রশ্ন ফাঁসের লক্ষণ & পরীক্ষা-অংশ নবায়ন & নি/\allowbreak{}মা & প্রতিটি বেঞ্চমার্ক সংস্করণে \\
R16 & গুরুত্বপূর্ণ জ্ঞান একজনের কাছে, কোনো সহকারী নেই & উত্তরাধিকার-পরিকল্পনা চালু & মা/\allowbreak{}মা & ত্রৈমাসিক \\
R17 & পুনর্নির্বাচনে দেখা যায় নির্বাচিত মডেল বিকল্পের চেয়ে ৫ পয়েন্টের বেশি পিছিয়ে & ADR-এর মাধ্যমে পুনর্নির্বাচন & নি/\allowbreak{}মা & বার্ষিক \\
R18 & বন্যা বা ঘূর্ণিঝড়ের সতর্কতা; ঘন ঘন বিদ্যুৎ-বিভ্রাট & পুনরুদ্ধার-সাইটে স্থানান্তর; অফলাইন পরিচালনা & মা/\allowbreak{}মা & প্রতি বর্ষার আগে ও মহড়ার পরে \\
\bottomrule
\end{tabularx}
\par\vspace{5pt}{\footnotesize\itshape অবশিষ্ট মাত্রা হলো সারণি ৫৫-এর প্রশমনের পর সম্ভাবনা ও প্রভাব: নি = নিম্ন, মা = মাঝারি, উ = উচ্চ। দায়িত্বপ্রাপ্তরা সারণি ৫৫-এ উল্লিখিত।}
\par\vspace{3pt}{\scriptsize\color{gray} উৎস: BAPIL কারিগরি ব্লুপ্রিন্ট, সংস্করণ ১.০ (২০২৬)। CC BY 4.0।}
\end{minipage}
\end{document}
